The creation of a bibliography is done in several steps.
Listing~\ref{lst:latex:bibliography:template} shows the basic structure.

\begin{compactenum}
\item First, the package~\ctanpkg{biblatex} must be loaded in the preamble.
\item Then, in the preamble with \lstinline[language={[LaTeX]{TeX}}]+\addbibresource+
      the literature databases are loaded. The template uses a single
      file~\texttt{literature.bib} by default. For extensive works or complex literature
      studies, it is also possible to combine the bibliography from several
      \texttt{.bib} files.
      It is advisable to store these files in a separate subdirectory.
      Each file with a corresponding
      \lstinline[language={[LaTeX]{TeX}}]+\addbibresource+-command can be
      loaded.
\item The actual bibliography is loaded within the document with
      \lstinline[language={[LaTeX]{TeX}}]+\printbibliography+
      is output.
\end{compactenum}

\lstinputlisting[language={[LaTeX]{TeX}},style=block,caption={Grundgerüst einer Bibliografie},label={lst:latex:bibliography:template}]{code/latex_bibliography_template.en.tex}

By default, the bibliography does not appear in the table of contents.
Listing~\ref{lst:latex:bibliography:toc} contains the necessary options to make
the bibliography appear in the table of contents.

\lstinputlisting[language={[LaTeX]{TeX}},style=block,caption={list bibliography in table of contents},label={lst:latex:bibliography:toc}]{code/latex_bibliography_toc.en.tex}

Both the package and the above commands can be customised with a variety of
other options. For this, please refer to the corresponding documentation of
\cite{pkg_biblatex}.

The source reference is done with the \lstinline[language={[LaTeX]{TeX}}]+\cite+
command. The keyword for the corresponds to the keyword assigned in the
\hologo{BibTeX} file. In addition, there are a large number of other citation
commands that are intended for different applications. The following paragraph
lists some variants by way of example. The \LaTeX code for this is in
Listing~\ref{lst:latex:bibliography:cite}. A detailed description of each variant
would go beyond the scope of this manual. For this, please refer to
\cite{pkg_biblatex,rees_biblatex_cheat}.

This documentation is published under \cite{urz_vorlage_manual}
published. The source reference can also be made by footnote
footnote.\footcite{urz_vorlage_manual} It is also possible to explicitly cite
cite only the authors: \citeauthor{urz_vorlage_manual} or
or just the title: \citetitle{urz_vorlage_manual}.
Even complete references in the body text are possible:
\fullcite{urz_vorlage_manual} or also in the footnote.
footnote.


\lstinputlisting[language={[LaTeX]{TeX}},style=block,caption={literature references},label={lst:latex:bibliography:cite},morekeywords={footcite,citeauthor,citetitle,fullcite,footfullcite}]{code/latex_bibliography_cite.en.tex}
